\originalchapter{光子关联与量子干涉}
\originalsection{光电探测与二阶关联}
理想光电探测的单计数率正比于$\langle E^{(-)}E^{(+)}\rangle$, 两次计数率正比于正规序$\langle E_1^{(-)}E_2^{(-)}E_2^{(+)}E_1^{(+)}\rangle$. 单模同点零延迟时,
\[g^{(2)}(0)=\frac{\langle a^{\dagger2}a^2\rangle}{\langle a^\dagger a\rangle^2}=1+\frac{\operatorname{Var}n-\bar n}{\bar n^2}.\]
相干态为1, 单模热态为2, 数态为$1-1/n$. 对由非负经典强度分布描述的场, $\langle I^2\rangle\ge\langle I\rangle^2$, 故$g^{(2)}(0)\ge1$; 低于1是非经典性的充分证据. 一个原子发射后必处于基态, 再激发需要时间, 所以出现反聚束; 探测器暗计数和有限时间分辨率会把理想凹陷填高.

Hanbury Brown--Twiss装置把场分成两路并比较计数, 不要求把“一个光子劈成两半”. 分束器后的联合正规序相关给同一输入的$g^{(2)}$, 真空输入的作用在保持对易关系中体现. 对热Gaussian场用Gaussian四阶矩分解得$g^{(2)}(\tau)=1+|g^{(1)}(\tau)|^2$, 单模假定及探测平均不能省略.

\originalsection{分束器与Hong--Ou--Mandel效应}
取$c=(a+b)/\sqrt2,d=(a-b)/\sqrt2$, 这是可吸收输出相位的分束器约定, 并保持$[c,c^\dagger]=[d,d^\dagger]=1$. 输入$\ket{1,1}=a^\dagger b^\dagger\ket{0,0}$变成$[(c^\dagger)^2-(d^\dagger)^2]\ket0/2=(\ket{2,0}-\ket{0,2})/\sqrt2$. 一路一个光子的振幅相消.

相消要求两个光子在偏振、谱形和时间上不可区分. 对归一化波包$f,g$, 重合概率为$P_{11}=[1-|\langle f|g\rangle|^2]/2$. 完全正交时回到经典$1/2$, 完全重叠时为零. 此式可由把$g$分成沿$f$和正交部分再分别演化得到. 干涉的是交换后的多粒子概率振幅, 不等于两束经典场的平均电场相消.

\originalsection{相位测量与损耗限制}
对Mach--Zehnder干涉仪, Schwinger算符$J_z=(n_a-n_b)/2,J_x=(a^\dagger b+b^\dagger a)/2,J_y=(a^\dagger b-b^\dagger a)/(2\ii)$满足角动量代数. 两分束器夹相移把输出布居差变成$J_z\cos\phi+J_x\sin\phi$. 在$N$个独立光子的最佳工作点, 误差传播$\delta\phi=\Delta J_z/|\partial_\phi\langle J_z\rangle|$给$1/\sqrt N$.

NOON态$(\ket{N,0}+\ket{0,N})/\sqrt2$积累相位$N\phi$, 理想局域灵敏度为$1/N$. 但一旦丢失一个光子, 环境便能记录哪一路的信息; 两臂均有透过率$\eta$时, 无丢失权重为$\eta^N$. 因此不能只写Heisenberg标度而不说明资源计数、先验相位范围和损耗. 固定非零损耗下大$N$的优势受限. 压缩态能在现实效率下改善常数, 也必须计入泵浦及非真空光子的资源.
